\documentclass[aspectratio=169,10pt,english]{beamer}

\input{../../../_preambulo_beamer.tex}
\input{../../../_comandos_beamer.tex}

\selectlanguage{english}

\title{MA1001B - Statistical Modeling for Decision Making}
\subtitle{Lesson 32: F Distribution and Comparing Two Variances}
\author{}
\institute{Tecnol\'ogico de Monterrey}
\date{}

\begin{document}

\begin{frame}[plain]
  \vspace{1.2cm}
  {\Large\bfseries MA1001B - Statistical Modeling for Decision Making\par}
  \vspace{0.7cm}
  {\large Lesson 32: F Distribution and Comparing Two Variances\par}
  \vspace{0.7cm}
  {\color{TecRojo}\rule{\textwidth}{1.2pt}}
  \vspace{0.8cm}

  MA1001B Faculty\\[0.3cm]
  Term\\[0.3cm]
  Tecnol\'ogico de Monterrey
\end{frame}

\begin{frame}{The Two-Variance Question}
  \begin{alertblock}{Guiding question}
    How can two independent sample variances be compared through an $F$
    ratio and an upper-tail probability?
  \end{alertblock}

  \vspace{0.6em}
  By the end of this lesson, you can:
  \begin{itemize}
    \item compute $s_{1}^{2}$ and $s_{2}^{2}$ from two samples;
    \item form $F=s_{1}^{2}/s_{2}^{2}$ with a pair of degrees of freedom;
    \item read $p$ from \texttt{scipy.stats.f.sf};
    \item explain why the two-variance $F$ test is sensitive to non-normality.
  \end{itemize}

  \vfill
  \scriptsize All cycle times used today are synthetic. No real company data are used.
\end{frame}

\begin{frame}{Two Independent Packing Stations}
  \small
  \begin{center}
    \begin{tabular}{lrr}
      \toprule
      \textbf{Quantity} & \textbf{Station A} & \textbf{Station B} \\
      \midrule
      Sample size & $n_{1}=25$ & $n_{2}=25$ \\
      Degrees of freedom & 24 & 24 \\
      Sample variance & $s_{1}^{2}=44.205731$ & $s_{2}^{2}=12.024693$ \\
      \bottomrule
    \end{tabular}
  \end{center}

  \vspace{0.5em}
  \begin{exampleblock}{Hypotheses}
    $H_{0}:\sigma_{1}^{2}=\sigma_{2}^{2}$ versus $H_{1}:\sigma_{1}^{2}>\sigma_{2}^{2}$.
    The stations are independent. Means are not the target of this test.
  \end{exampleblock}
\end{frame}

\begin{frame}[fragile]{Step 1: Two Sample Variances}
  \begin{columns}[T]
    \begin{column}{0.42\textwidth}
      \small
      \begin{lstlisting}
s1_sq = np.var(a, ddof=1)
s2_sq = np.var(b, ddof=1)
      \end{lstlisting}

      \begin{block}{Verified output}
        $s_{1}^{2}=44.205731$\\
        $s_{2}^{2}=12.024693$\\
        $\mathrm{df}=(24,24)$
      \end{block}
    \end{column}
    \begin{column}{0.58\textwidth}
      \centering
      \includegraphics[width=\textwidth]{../figures/f_two_var_01_sample_variances.png}
      \vspace{0.2em}
      \scriptsize Station boxplots from Step 1
    \end{column}
  \end{columns}
\end{frame}

\begin{frame}{The $F$ Ratio and the Upper Tail}
  \small
  \[
    F=\frac{s_{1}^{2}}{s_{2}^{2}}=\frac{44.205731}{12.024693}=3.676246.
  \]
  Critical value and $p$-value:
  \[
    F_{0.05}(24,24)=1.983760,
  \]
  \[
    p=\texttt{f.sf}(3.676246,24,24)=0.001125.
  \]
  Because $p<\alpha=0.05$, $H_{0}$ is rejected: Station A is more variable
  in this synthetic snapshot.
\end{frame}

\begin{frame}[fragile]{Step 2: $F$ and \texttt{f.sf}}
  \begin{columns}[T]
    \begin{column}{0.40\textwidth}
      \small
      \begin{lstlisting}
F = s1_sq / s2_sq
p = stats.f.sf(F, 24, 24)
      \end{lstlisting}

      \begin{block}{Verified output}
        $F=3.676246$\\
        $p=0.001125$\\
        Reject $H_{0}$
      \end{block}
    \end{column}
    \begin{column}{0.60\textwidth}
      \centering
      \includegraphics[width=\textwidth]{../figures/f_two_var_02_f_test.png}
      \vspace{0.2em}
      \scriptsize $F$ upper tail from Step 2
    \end{column}
  \end{columns}
\end{frame}

\begin{frame}{Step 3: $F$ Is Sensitive to Non-Normality}
  \begin{columns}[T]
    \begin{column}{0.42\textwidth}
      \small
      Equal variances, 5000 tests, $\alpha=0.05$.

      \vspace{0.4em}
      Normal samples: Type I rate $0.052200$.

      \vspace{0.4em}
      $t(3)$ samples: Type I rate $0.202400$.

      \vspace{0.5em}
      \textbf{Limit:} heavy tails inflate false rejections even when the two
      variances are equal. The two-variance $F$ test is not robust.
    \end{column}
    \begin{column}{0.58\textwidth}
      \centering
      \includegraphics[width=\textwidth]{../figures/f_two_var_03_nonnormal_sensitivity.png}
      \vspace{0.2em}
      \scriptsize Type I rates from Step 3
    \end{column}
  \end{columns}
\end{frame}

\begin{frame}{Concept Check}
  \begin{enumerate}
    \item Compute $F$ from $s_{1}^{2}=44.205731$ and $s_{2}^{2}=12.024693$.
    \item What are the two degrees of freedom for $n_{1}=n_{2}=25$?
    \item Does $p=0.001125$ compare means or variances?
    \item Why is a Type I rate of $0.202400$ a problem at $\alpha=0.05$?
  \end{enumerate}

  \vspace{0.6em}
  \begin{block}{Connection to Lesson 33}
    The session can continue directly into $H_{0}$, $H_{1}$, Type I and Type
    II errors as a general testing framework.
  \end{block}

  \vfill
  \scriptsize Source: OpenStax, \textit{Introductory Business Statistics 2e},
  Chapter 12, Section 12.1. CC BY-NC-SA.
\end{frame}

\appendix



\end{document}
